\section{东北、燕云、河南：同一王朝的不同时间}

\subsection{上京方向：建国的资源仍留在东北}

金的早期中心在东北，并不因为后来中都成为都城而失去意义。按出虎水、会宁一带的河谷、草场、城寨和通向辽东、松花江流域的路线，是完颜部兴起时能够集结人马、交换物资和传递消息的条件。金朝的国号把这些条件政治化，却没有把它们变成可以无限抽取的资源。寒暑、道路、马匹和邻近部众的关系，都会限制一项征发可以走多远、多久返回。

东北地方的材料难以恢复。正史倾向于从太祖的军事行动写起，后来的地理志又容易将定制投回创业阶段。女真文字、汉文和契丹文字材料能够在少数场合补足官号、姓名和营建活动，但释读、断代和保存来历仍须逐一辨认。它们不支持把东北想象为没有城市、没有贸易、只有单一部众的“边缘”；也不支持把所有不同语言的题记解释为同一种政治认同。能确认的是，金最初的军政网络依赖这一区域，而进入华北后的统治必须持续把东北的人力、记忆和机构同新的中心联系起来。

猛安谋克在这条联系中具有特殊位置。它既与部众、军役和土地关系相关，也会随迁徙、任官和扩张改变具体面貌。把它译成一个固定的现代制度名称，往往掩盖它在东北与华北可能承担的不同任务。某些猛安谋克成员进入华北，既是统治的延伸，也会改变当地军役、居住和资源分配；当地居民则未必以同样方式被纳入。朝廷需要这样的组织维系军队，地方却要在驻军、耕作、市场和亲属关系之间承受它的日常后果。

\subsection{燕云：交接地而不是奖品}

燕云在宋金联盟中被当作目标，在辽亡后却迅速成为难题。其城镇、关隘和道路连接北方草原、东北与华北平原；谁占据它，不只获得象征性的旧疆，也要接管守备、仓储和地方关系。金把部分地区交与宋又收回的过程，说明政治承诺不足以安排地面控制。宋军的进入、金军的撤留、原辽官民的选择，都发生在消息不完全、粮草有限的条件下。

燕京后来被改造为金中都，这个转变不能掩盖先前的断裂。战争会破坏既有运输和市场，也会使某些军需、修建和交易集中起来。地方精英可能利用熟悉的文书和人际网络保存位置，也可能因依附旧朝而被排斥；商人和工匠可能随军队、市场和工程改变居处。我们没有足够连续的资料逐户追踪这些变化，因此不以“合作”或“抵抗”两词替代复杂选择。城址显示的营建层和墓志所记的迁居，只能让我们在具体地点看见局部的再组织。

燕云也是语言和礼仪交错的地方。辽、宋、金各自的官名、地名和外交称谓不会在易主当天消失。文书可被沿用、改写或重新解释；一个人拥有的官衔，也可能来自实授、遥授、追赠或不同政权的承认。碑志中的称谓因此需要放回成文情境，而不能直接作为身份或忠诚的证据。读者若只问谁“属于”燕云，会错过更紧迫的问题：谁能调兵、收粮、开市和裁决纠纷。

\subsection{河北与山东：平原上的供给和脆弱}

华北平原使大规模运输成为可能，也使战事能够沿道路与河道扩散。河北、山东的城镇、粮区和市场，对于金的南下与中都供给都很重要；这种重要性也让它们更容易承受驻军、征粮和反复易手。平原并不意味着没有阻碍。河道改道、堤防失修、雨雪和桥渡状况都能改变一支部队的速度；仓储在何处、车辆由谁提供、民夫何时脱离农事，同样决定一项计划是否可行。

州县文书通常保存得不完整，故我们不应把《金史》中的一项调粮令当作所有村落的经验。灾异记录也要谨慎。旱、蝗、水灾被写入帝纪，可能是地方确有严重困境，也可能被纳入君主修德与官员考课的政治语言；它不能自动折算成某一地区的死亡率或全国粮价。更可靠的做法是把灾害放回道路、赋役和迁徙：如果一段堤防失效，影响的是渡运、田地和仓粮；如果军队同时过境，家庭面对的风险就会改变。因果在此是相互叠加，不是自然灾害单独决定王朝命运。

山东沿海与内陆市场又提醒我们，金并非只有一条朝向南宋的战线。盐、海运、手工业和区域交换把地方同更远的网络联系起来，国家可以从中征课，也必须容忍大量难以完全控制的中间人。考古所见的瓷器、钱币和港口遗存可以显示物资的到达和使用，不能直接等于官方贸易路线或人口财富。它们最有价值之处，在于迫使我们把“北方”理解为不同的水陆网络，而非只是一片等待中枢征收的土地。

\subsection{河南：迁都之后的第二个中心}

一二一三年后金廷向南调整，开封及河南的地位陡然上升。选择南迁不等于抛弃北方时毫无代价的理性决定，也不等于只因君主怯懦。中都面临威胁，朝廷需要一个仍可组织官员、粮食、军队和外交的中心；河南有城市和农业资源，也已承受南宋边界和过境军队的压力。迁入者带来的官署、军人、工匠和需求会改变市场与房舍，新的中枢则需要重新建立仓储、邮传和地方合作。

中都失守后，旧中心并未从人们的路线图中抹去。逃离者、留下者、后来进入者各自把道路记成不同的危险与机会。河南也不能在一纸迁都令后立刻成为稳定腹地。黄河、淮北方向的防御、税粮转运、军队开支和地方灾害会不断把边区压力推向新的都城。金廷越需集中资源，地方越可能感到征发和官吏的存在；地方一旦难以供给，前线又会反过来威胁中枢。

一二三二年汴京受围时，这种循环达到极端。围城不是单纯的军事事件：它压缩了粮食、燃料、疾病、价格和信息，令原本分散在区域间的脆弱一同进入城市。史书中关于守城的忠勇和责任，能说明政治如何理解危机；它们无法替每个居民保存经验。我们可以说，围困使金廷失去一个重要的行政与供给节点；不能把城内所有人写成同一种受害者或同一种支持者。

\subsection{从区域差异理解金的社会}

金朝的空间从东北河谷、燕云城市到华北平原、河南中枢和淮河边防，各处面对的征服、行政与财政时间不同。中都的政策可能先变成官署和工程，边寨先感到的却是粮车和调兵；东北的军政传统可以被调用，河北的家庭则可能先遇到新税目或驻军；河南在末期成为中枢时，还要承担来自北方与南方的双重压力。以一条“金朝制度”概括这些差异，会让国家看似强大而社会消失。

反过来，区域差异不意味着各地互不相干。道路、河流、官员调动、市场和战争不断把它们接在一起。金的统治能力正在于让这些连接能够传递命令、粮食和人力；它的脆弱也在于一地的失收、城陷或道路中断会传到另一地。晚金面对蒙古攻势时，真正被侵蚀的不是抽象领土，而是维系区域间关系的通道与信任。
